\documentclass[11pt]{article}

\usepackage{amsmath,amssymb,amsthm}
\usepackage[margin=1.2in]{geometry}
\usepackage{xcolor}
\usepackage[colorlinks=true,linkcolor=blue,citecolor=blue]{hyperref}

\newtheorem{theorem}{Theorem}[section]
\newtheorem{lemma}[theorem]{Lemma}
\newtheorem{proposition}[theorem]{Proposition}
\newtheorem{corollary}[theorem]{Corollary}
\theoremstyle{definition}
\newtheorem{definition}[theorem]{Definition}
\newtheorem{example}[theorem]{Example}
\newtheorem{question}[theorem]{Question}
\theoremstyle{remark}
\newtheorem{remark}[theorem]{Remark}

\newcommand{\F}{\mathbb{F}}
\newcommand{\Z}{\mathbb{Z}}
\newcommand{\Q}{\mathbb{Q}}
\newcommand{\Clo}{\mathrm{Clo}}
\newcommand{\Fbar}{\overline{\F}_2}

\title{Single generators of polynomial rings under composition:\\
a characteristic-two dichotomy\\[1ex] {\large\normalfont --- draft ---}}
\author{John Thompson}
\date{June 2026}

\begin{document}
\maketitle

\begin{abstract}
In Boolean logic a single operation, NAND, suffices to express every
Boolean function.  We ask the analogous question for polynomial rings:
is there a single binary polynomial $p$ such that every polynomial in
two variables can be written as an iterated composition of $p$
(starting from the variables and constants)?  Over $\Z$ the answer is
no: for every $p \in \Z[x,y]$, at least one of $x+y$, $xy$ is
inexpressible.  More generally, there is a single polynomial
generating all of $R[x,y]$ under composition (with variables and
constants) if and only if $2$ is invertible in $R$: when it is, the
single operation $p(x,y) = x^2 - y$ generates everything; when it is
not, generation fails through a residue field of characteristic $2$.
The main work is the negative result
over fields of characteristic $2$, proved by an infinite descent
whose decreasing quantity is, informally, the number of times a
would-be counterexample factors through the squaring map.  All
results have been checked formally.
\end{abstract}

\section{Introduction}

\subsection{A Sheffer stroke for polynomial rings}

Every Boolean function can be built from a single connective: NAND
(the \emph{Sheffer stroke}) \cite{Sheffer,Post}.  One-element generating sets of this kind
appear throughout algebra and logic, and it is natural to ask for an
analogue one level up, with the two-element Boolean domain replaced by
a ring of polynomials.

So fix the ring $\Z$ and consider polynomials in two variables.  We
are allowed to start from the raw materials
\[
  x, \qquad y, \qquad \text{and every constant } c \in \Z,
\]
and we are given \emph{one} binary polynomial $p(x,y)$ as our only
``connective'': whenever $f$ and $g$ have been built, we may build
$p(f,g)$.  Write $\Clo(p)$ for the set of all polynomials obtainable
this way (a precise definition is in Section~\ref{sec:prelim};
Section~\ref{sec:precision} explains how this relates to the standard
notion of a \emph{clone} in universal algebra; no background from
that subject is needed).  Say that $p$ \emph{generates} $\Z[x,y]$ if
$\Clo(p) = \Z[x,y]$.

Could such a $p$ exist?  Nothing obvious rules it out; a single
polynomial has plenty of room to encode both an ``additive'' and a
``multiplicative'' mode, depending on what is substituted into it.
And since every polynomial is built from variables and constants by
$+$ and $\times$, generation is equivalent to just two
memberships (Lemma~\ref{lem:criterion}):
\[
  \Clo(p) = \Z[x,y] \iff x + y \in \Clo(p) \ \text{ and } \
  xy \in \Clo(p).
\]
Our first result is that no single polynomial achieves even this.

\begin{theorem}\label{thm:Z}
For every $p \in \Z[x,y]$, at least one of $x+y$, $xy$ lies outside
$\Clo(p)$.  In particular, no binary polynomial generates $\Z[x,y]$.
\end{theorem}

\subsection{Results}\label{sec:results}

Theorem~\ref{thm:Z} is the case $R = \Z$ of a dichotomy over
arbitrary commutative rings, which is the main result of this paper:

\begin{theorem}[the dichotomy]\label{thm:dichotomy}
Let $R$ be a commutative ring.  There is a single polynomial
generating all of $R[x,y]$ under composition (with variables and
constants) if and only if $2$ is a unit of $R$.
\end{theorem}

Half of this is easy: over any commutative ring in which $2$
is invertible, the operation
\[
  p(x,y) \;=\; x^2 - y
\]
generates the full polynomial ring --- a dozen lines of explicit
algebra, presented in
Section~\ref{sec:converse}, valid over $\Q$, $\mathbb{R}$,
$\Z[\tfrac12]$, and $\F_p$ for every odd prime $p$.  (Over $\F_3$,
for example, iterating $x^2 - y$ from the variables and constants
produces \emph{every} polynomial in $\F_3[x,y]$.)

\begin{theorem}\label{thm:converse}
Let $R$ be a commutative ring with $\tfrac12 \in R$.  Then
$q(x,y) := x^2 - y$ generates the full polynomial ring:
$\Clo(x^2 - y) = R[x,y]$.
\end{theorem}

The same derivation, run over $\Z$, fails exactly at the steps that
divide by $2$; Theorem~\ref{thm:Z} says this failure is no accident
of one derivation.

The negative half starts from the observation
that generation survives reduction of coefficients: a ring
homomorphism $R \to S$ carries any derivation of $x+y$ over $R$ ---
by a \emph{derivation} we mean a formula building $x+y$ from the
variables and constants by repeated application of $p$ --- to one
over $S$.  Reducing mod $2$, Theorem~\ref{thm:Z} therefore follows
from the corresponding statement over the two-element field
$\F_2 = \Z/2$:

\begin{theorem}\label{thm:F2}
For every $q \in \F_2[x,y]$,
\[
  x + y \notin \Clo(q) \qquad\text{or}\qquad xy \notin \Clo(q).
\]
\end{theorem}

Two remarks on the statement.  First, polynomials here are
\emph{formal} objects, not functions: over $\F_2$ infinitely many
polynomials induce the same function on $\F_2^2$, and
Theorem~\ref{thm:F2} concerns the polynomials themselves, so it does
not follow from (and does not contradict) the classical theory of
functionally complete operations on finite sets;
Section~\ref{sec:precision} dwells on the distinction.  Second ---
and this matters for the reduction --- there is \emph{no hypothesis
on $q$ whatsoever}.  Reduction mod $2$ can collapse degree
($2x^3 + xy \mapsto xy$), so to deduce Theorem~\ref{thm:Z} one needs
the $\F_2$-statement for arbitrary $q$, including the
low-degree and degenerate ones.

Passing to $\F_2$ is not merely a convenience: characteristic $2$ is
where the obstruction lies, and not only over the prime
field.  The descent behind Theorem~\ref{thm:F2} generalizes to every
\emph{perfect} field of characteristic $2$ (one in which every
element has a square root --- all finite fields $\F_{2^k}$ and all
algebraically closed fields, for instance), at the cost of a single
coefficient twist explained in Section~\ref{sec:tothedichotomy}, and
the disjunctive conclusion of Theorem~\ref{thm:F2} holds verbatim
there; an arbitrary field of characteristic $2$ then embeds into its
perfection, and non-membership in the clone descends along
coefficient embeddings.  Hence:

\begin{theorem}\label{thm:char2}
For every field $k$ of characteristic $2$ and every $q \in k[x,y]$,
at least one of $x+y$, $xy$ lies outside $\Clo(q)$.  In particular,
over a field $k$ of characteristic $2$ no single binary polynomial
generates $k[x,y]$.
\end{theorem}

Finally, if $2$ is not a unit of $R$, it lies in a maximal ideal; a
generator would push forward along the quotient map to the
residue field, which has characteristic $2$.
Section~\ref{sec:tothedichotomy} assembles these reductions into a
proof of Theorem~\ref{thm:dichotomy}.

\subsection{What exactly is claimed}\label{sec:precision}

Because the objects here are formal polynomials rather than
functions, some care is needed about what
Theorems~\ref{thm:Z}--\ref{thm:char2} do and do not assert; this
subsection collects the necessary precision.

In universal algebra, a \emph{clone} on a set $A$ is a family of
finitary operations $A^n \to A$ containing all coordinate projections
and closed under composition; the classical theory --- Post's lattice
and its descendants --- studies such families of \emph{functions}
\cite{Post,Szendrei}.  The elements of $\Clo(p)$ are instead formal
polynomials, not maps.  In clone-theoretic terms: the binary terms in
the language of commutative rings with constant symbols form the
binary part of an abstract clone (equivalently, of a Lawvere theory
\cite{Lawvere}), canonically identified with $R[x,y]$ under
substitution, and our question asks when the single term $p$
generates it.  (Restricting to the binary part loses nothing: any
term built from $p$, constants, and two variables is itself a
bivariate polynomial, whatever arities its subterms pass through.)
This abstract object is represented faithfully by operations on $R$
when $R$ is an infinite integral domain, and is finer than its image
otherwise; it should not be confused with the clone of polynomial
\emph{operations} of the ring $R$ in the sense of the
universal-algebra literature \cite{KaarliPixley,Szendrei}, nor with
the full clone of all finitary operations on the set $R$.  A reader
who prefers to avoid clones altogether can take the low-tech view of
Definition~\ref{def:clone} at face value: $R[x,y]$ carries the single
binary operation $(f,g) \mapsto p(f,g)$, and $\Clo(p)$ is the
subalgebra of this magma generated by $\{x, y\} \cup R$.

Over an infinite integral domain a polynomial is determined by the
function it induces, so over $\Z$ or $\Q$ our statements can equally
be read as statements about operations (on the infinite set $\Z$ or
$\Q$, where the classical finite-domain theory is silent).  Over
finite fields the formal question is strictly finer than the
functional one, and the contrast is sharpest there: by a theorem of
Webb \cite{Webb}, a single binary operation generates every
operation on any finite set --- in particular every polynomial
function over $\F_2$ is a composition of a single binary polynomial
function --- while Theorem~\ref{thm:F2} shows that no single
polynomial generates $\F_2[x,y]$ formally.  The formal statement
therefore has no functional counterpart.

\begin{example}\label{ex:nand}
Concretely, $n(x,y) := 1 + xy \in \F_2[x,y]$ induces the binary NAND
operation on $\F_2$, and every Boolean function is a composition of
NAND; nevertheless $x + y \notin \Clo(n)$ by Theorem~\ref{thm:F2},
since $Dn = 1 \neq 0$ for the mixed-derivative operator $D$ of
Section~\ref{sec:dichotomy}.
\end{example}

Two further cautions about nearby terminology.  First, we admit all
constants as generators; in the function-clone world this corresponds
to \emph{completeness with constants} in the classical functional
sense.  (The constant-free variant is also natural; we do not pursue
it here.)  Second, the phrase \emph{polynomial completeness} has an
established meaning in universal algebra --- that every
congruence-compatible function on an algebra is induced by a
polynomial, as in the monograph \cite{KaarliPixley} --- unrelated to
the generation property studied here.

\subsection{Relation to prior work}\label{sec:priorwork}

Generation from a single binary operation has a long history in the
function setting.  Webb \cite{Webb} showed that on any finite set a
single binary operation generates all finitary operations under
composition; since every operation on a finite field is induced by a
polynomial, a single binary polynomial function generates all
polynomial functions over $\F_q$ for every prime power $q$ ---
including in characteristic $2$, where the formal generation problem
has a negative answer (Theorem~\ref{thm:char2}); the two settings
were compared in Section~\ref{sec:precision}.  On infinite domains,
Sierpi\'nski \cite{Sierpinski} proved that every finitary operation
is a composition of binary ones, and set-theoretic encoding
arguments produce a single binary operation on $\Z$ from which
addition and multiplication are both recovered as terms
\cite{HamkinsMO}; for the clone theory of infinite domains see the
survey of Goldstern and Pinsker \cite{GoldsternPinsker}.  The
operations these constructions produce are far from polynomial, and
the surrounding theory --- Sheffer operations and the functional
completeness criteria recalled in Section~\ref{sec:context} ---
concerns clones of arbitrary functions rather than formal terms.  In
a different direction, Odrzywo{\l}ek \cite{Odrzywolek} recently
exhibited a single binary operation built from $\exp$ and $\log$
that generates the elementary functions of analysis by composition
with the constant $1$ --- the same question for functions of a real
variable rather than formal polynomials, with no negative half.

The closest prior work we have found concerns generation of integer
polynomials when addition is available alongside composition.
Aichinger \cite{Aichinger2013} characterized the univariate integer
polynomials obtainable from $x$ and $x^2$ (or from $1$, $x$, and
$x^3$) under addition, subtraction, and composition, and Aichinger
and Kreinecker \cite{AichingerKreinecker} described the subnearrings
of $(\Z[x], +, \circ)$ generated by the subsets of
$\{1, x, x^2, x^3\}$, proving non-membership results such as
$2x^{10} \in \langle x^2 \rangle$ but
$x^{10} \notin \langle x^2 \rangle$ --- a $2$-divisibility
obstruction to composition generation over $\Z$ that echoes, in a
different signature, the role parity plays here.  Those results are
univariate, are confined to $\Z$, and retain addition as a
primitive; the present setting removes addition entirely, asks for a
single binary generator of the full bivariate polynomial ring over
an arbitrary commutative ring, and identifies the invertibility of
$2$ as the exact boundary.  We have not found the composition-only
single-generator question for formal polynomial rings in the
literature; Section~\ref{sec:context} returns to the adjacent
literatures on functional completeness and on composition structures
of polynomial rings.

\subsection{The shape of the negative proof, in words}\label{sec:shape}

The substantial result is Theorem~\ref{thm:F2}, and the rest of this
introduction previews its proof at leisure; the new idea is the
descent in step 3, and we will take it slowly again when it appears
for real in Section~\ref{sec:descent}.

\medskip
\noindent\emph{Step 1: a derivative splits the problem.}
Characteristic $2$ kills second derivatives in each single variable
($\partial_x^2 = \partial_y^2 = 0$) but not the mixed one, so the
operator $D := \partial_x \partial_y$ carries unusual weight (all
derivatives in this paper are formal partial derivatives of
polynomials).  We split according to $Dq$.

If $Dq = 0$ --- that is, no monomial of $q$ has both exponents odd ---
then a chain-rule identity shows the property ``$Dg = 0$'' is
preserved by every substitution into $q$.  The variables and constants
have it; $xy$ does not ($D(xy) = 1$).  So $xy \notin \Clo(q)$, and
this case is done.

If $Dq \neq 0$ we will show $x + y \notin \Clo(q)$; this case takes
the rest of the argument.

\medskip
\noindent\emph{Step 2: from clones to curves.}
To exclude $x+y$ we use an invariant of a softer kind: a rigidity
property of substitution.  We will show that, when $Dq \neq 0$, a
substitution $q(\alpha,\beta)$ can only land on a nonconstant
polynomial in $x+y$ if $\alpha$ and $\beta$ are themselves
polynomials in $x+y$.  A short induction over the clone
(Proposition~\ref{prop:master}) shows this rigidity implies
$x+y \notin \Clo(q)$.

The rigidity property, in turn, is a statement about the \emph{level
sets} of $q$ --- the loci $\{q = c\}$, $c$ a constant.
Restricting everything to a generic line of the pencil
$\{x + y = \text{const}\}$ of parallel lines converts it into the
following concrete
geometric assertion: there is no pair of one-variable polynomials
$\alpha(t), \beta(t)$, not both constant, with
\[
  q\bigl(\alpha(t), \beta(t)\bigr) = c
\]
identically in $t$, where $c$ is a ``generic'' constant --- precisely,
a constant transcendental over $\F_2$ (satisfying no nonzero
polynomial relation with $\F_2$-coefficients), living in a large
algebraically closed field $K$ of characteristic $2$; the setting is
fixed precisely in Section~\ref{sec:bridge}.  Such a pair
$(\alpha, \beta)$ would be a polynomial-parametrized curve lying
inside a single generic level set of $q$; we call it a
\emph{witness} (Definition~\ref{def:witness}).  The task is to show witnesses do not exist when
$Dq \neq 0$.

\medskip
\noindent\emph{Step 3: the descent.}
Here is the idea in one paragraph.  Differentiate the witness equation
$q(\alpha,\beta) = c$ with respect to $t$.  In characteristic $2$,
derivatives interact with squares in extreme ways (squares have
derivative zero; conversely, over a perfect field, anything with
derivative zero \emph{is} a square), and the derivative calculus along
the witness forces a great deal of square structure.  The endgame is a
case split.  Either some polynomial relation with $\F_2$-coefficients holds
along the witness --- and any such relation pins
the witness to algebraic data, contradicting the genericity of the
level $c$ --- or no such relation holds, in which case both $\alpha$
and $\beta$ turn
out to be squares of polynomials of half the degree,
$\alpha = \alpha_1^2$, $\beta = \beta_1^2$, and because the
coefficients of $q$ lie in $\F_2$ (where squaring fixes everything),
\[
  q(\alpha_1, \beta_1)^2 = q(\alpha_1^2, \beta_1^2) = c,
\]
so $(\alpha_1, \beta_1)$ is again a witness, for the \emph{same} $q$,
at the level $\sqrt{c}$, with \emph{half the degree}.  Repeat.  The
degree cannot halve forever, so eventually the first branch occurs.

The induction, note, is not on $q$ at all: $q$ never changes.  It is
on the witness, and the decreasing quantity is the number of times the
witness can still factor through the Frobenius (squaring) map --- what
one might call its \emph{Frobenius depth}.  This inversion ---
inducting on the counterexample's Frobenius depth rather than on
any complexity measure of $q$ --- is the main idea of the proof; our
attempts at induction on degree-like measures of
$q$ had all failed.  Degree of $q$ is a poor induction variable here because
substitution composes degrees in uncontrollable ways (note e.g.\ that
$q$ may reproduce low-degree polynomials from high-degree inputs);
Frobenius depth of the witness is immune to this.

\medskip
Because the argument has an unusual density of characteristic-$2$
bookkeeping, we have verified all results formally; a short note on
this appears in Section~\ref{sec:lean}, and the text otherwise
presents proofs as ordinary mathematics.  Specifically, every result
is checked down to the proof kernel in Lean, and the headline theorems
are additionally verified by an adversarial judge
(\texttt{leanprover/comparator}) against a standalone statement file
of roughly ninety lines, with every proof replayed through two
independent kernels --- so the many edge cases, where an argument this
dense is most fragile, are covered by the verification.

\section{Definitions and conventions}\label{sec:prelim}

\begin{definition}\label{def:clone}
Let $R$ be a commutative ring, $R[x,y]$ the ring of formal
polynomials in two variables, and $p \in R[x,y]$.  The set
$\Clo(p) \subseteq R[x,y]$ is the smallest
set of polynomials such that
\begin{enumerate}
  \item $x \in \Clo(p)$, $y \in \Clo(p)$, and $c \in \Clo(p)$ for
        every constant $c \in R$;
  \item if $f, g \in \Clo(p)$ then $p(f,g) \in \Clo(p)$,
\end{enumerate}
where $p(f,g)$ denotes substitution of $f$ for $x$ and $g$ for $y$ in
$p$.  We say $p$ \emph{generates} $R[x,y]$ if $\Clo(p) = R[x,y]$.
\end{definition}

Thus $\Clo(p)$ is the set of polynomials expressible by a formula in
the variables and constants whose only operation symbol is $p$.

\begin{lemma}\label{lem:criterion}
$\Clo(p) = R[x,y]$ if and only if $x + y \in \Clo(p)$ and
$xy \in \Clo(p)$.
\end{lemma}

\begin{proof}
One direction is trivial.  For the other, first note that substituting
members of $\Clo(p)$ into a \emph{member} of $\Clo(p)$ again lands in
$\Clo(p)$ (induction on the formula).  So if $x+y$ and $xy$ are
members, $\Clo(p)$ is closed under the derived binary operations
$(f,g) \mapsto f + g$ and $(f,g) \mapsto fg$; since it contains the
variables and constants, it contains every polynomial.
\end{proof}

\begin{lemma}[transfer]\label{lem:transfer}
Let $\varphi : R \to S$ be a ring homomorphism and
$\pi : R[x,y] \to S[x,y]$ the induced map on coefficients.  Then
$\pi\bigl(\Clo(p)\bigr) \subseteq \Clo(\pi p)$ for every
$p \in R[x,y]$.
\end{lemma}

\begin{proof}
$\pi$ fixes the variables, sends constants to constants, and commutes
with substitution; induct over the formula.
\end{proof}

In particular, if both $x+y, xy \in \Clo(p)$ over $\Z$, then both lie
in $\Clo(\pi p)$ over $\F_2$.  So Theorem~\ref{thm:F2} (for
\emph{arbitrary} $q = \pi p$, degenerate cases included) implies
Theorem~\ref{thm:Z}, and likewise its analogue over any ring with a
homomorphism to $\F_2$ ($\Z/2^k$, $\F_2$-algebras, \dots).
Sections~\ref{sec:dichotomy} and~\ref{sec:descent} take place
entirely over $\F_2$; general coefficients return in
Section~\ref{sec:tothedichotomy}.

\medskip

We will constantly use three elementary facts about the formal partial
derivatives $\partial_x, \partial_y$ on $\F_2[x,y]$, all special to
characteristic $2$:

\begin{enumerate}
  \item \emph{Squares are derivative-free}:
        $\partial(f^2) = 2f\,\partial f = 0$.
  \item \emph{Repeated differentiation in one variable dies}:
        $\partial_x^2 = \partial_y^2 = 0$ (a monomial surviving
        $\partial_x$ has had its odd $x$-exponent made even), while the
        mixed operator $D := \partial_x\partial_y$ does not.
  \item \emph{Frobenius}: squaring is a ring homomorphism,
        $(f+g)^2 = f^2 + g^2$; and a polynomial whose exponents are
        all even is a square (halve the exponents and use that every
        coefficient, being $0$ or $1$, is its own square root).
\end{enumerate}

One further convention.  Degrees of polynomials in one variable are
taken with the convention that constants --- including the zero
polynomial --- have degree $0$; under this convention
$\deg(f^2) = 2\deg f$ holds without exception, which keeps the
bookkeeping of the descent in Section~\ref{sec:descent} honest.

\section{The positive half: $x^2 - y$}
\label{sec:converse}

We begin with the positive half of the dichotomy, which is
self-contained and shows concretely what the negative half must rule
out.

\begin{proof}[Proof of Theorem~\ref{thm:converse}]
By Lemma~\ref{lem:criterion} it suffices to reach $x+y$ and $xy$.
All steps below are instances of one identity:
\begin{equation}\label{eq:master-id}
  q\bigl(\alpha,\, q(\alpha', \gamma)\bigr)
  = \alpha^2 - \alpha'^2 + \gamma .
\end{equation}

\emph{Translations.}  With constants
$\alpha = \frac{d+1}{2}$, $\alpha' = \frac{d-1}{2}$,
identity \eqref{eq:master-id} sends $\gamma \mapsto \gamma + d$, for
any prescribed $d \in R$.  (This is the only step that uses
$\tfrac12$.)  In particular $x + d, y + d \in \Clo(q)$ for all $d$.

\emph{Slanted lines.}  $q(c, y) = c^2 - y$, then
$q(x, c^2 - y) = x^2 + y - c^2$, then
\[
  q\bigl(x + d,\; x^2 + y - c^2\bigr) = 2dx - y + (d^2 + c^2)
  \;\in\; \Clo(q),
\]
a line of slope $2d$ for every $d$, with constant adjustable by
translations.

\emph{Linear span step.}  For any $\alpha \in \Clo(q)$ and constant $k$,
apply \eqref{eq:master-id} with outer argument $\alpha + k$ and inner
argument $\alpha$: the increment is
\[
  \gamma \;\longmapsto\; (\alpha + k)^2 - \alpha^2 + \gamma
  \;=\; \gamma + 2k\,\alpha + k^2 .
\]
Increments along slanted lines of two distinct slopes span all linear
polynomials, so starting from $\gamma = 0$: \emph{every linear
polynomial lies in $\Clo(q)$}, in particular $x + y$ and
$x + \tfrac{y}{2}$.

\emph{Scaled squares.}  The same increment step with $\alpha = y^2 + \ell$
(in the clone via $q(y, -\ell)$) adds $\lambda y^2 + \text{const}$ for
any $\lambda$; from $\gamma = x^2 = q(x, 0)$ we get
$x^2 + \tfrac14 y^2 \in \Clo(q)$.

\emph{Finish.}
\[
  q\Bigl(x + \tfrac{y}{2},\; x^2 + \tfrac14 y^2\Bigr)
  = \Bigl(x + \tfrac{y}{2}\Bigr)^2 - x^2 - \tfrac14 y^2
  = xy . \qedhere
\]
\end{proof}

\begin{remark}
All constants used lie in $\Z[\tfrac12]$, so the derivation runs
verbatim over every ring in which $2$ is invertible --- in particular
over $\Q$, $\mathbb{R}$, $\Z[\tfrac12]$, and $\F_p$ for odd $p$; over
$\F_3$, the single operation $x^2 - y$ generates the full polynomial
ring $\F_3[x,y]$.  Over $\Z$ the translation step requires
realizing $d = c^2 - e^2$ with $c, e \in \Z$, which is possible only
for $d \not\equiv 2 \pmod 4$; the increments available are then too
even to span, and the derivation breaks down --- as, by
Theorem~\ref{thm:Z}, every derivation must.
\end{remark}

\begin{remark}
Theorem~\ref{thm:converse} produces a generator of degree
$2$.  We do not know the full picture for higher degrees; see
Section~\ref{sec:open}.
\end{remark}

\section{The negative half over $\F_2$}\label{sec:dichotomy}

We turn to the negative half.  This section and the next prove
Theorem~\ref{thm:F2}.  Fix $q \in \F_2[x,y]$.

\subsection{Parity decomposition and the derivative dichotomy}

Group the monomials $x^i y^j$ of $q$ into four classes by the parities
of $(i,j)$.  Each class is an even-exponent polynomial times one of
$1, x, y, xy$, and even-exponent polynomials are squares (fact 3
of Section~\ref{sec:prelim}).  Hence $q$ can be written, uniquely, as
\begin{equation}\label{eq:parity}
  q \;=\; A_0^2 \;+\; x A_1^2 \;+\; y A_2^2 \;+\; xy\,A_3^2,
  \qquad A_i \in \F_2[x,y].
\end{equation}
Differentiating \eqref{eq:parity} and discarding the squares
(fact 1), only the prefixes $x$, $y$, $xy$ contribute:
\begin{equation}\label{eq:partials}
  \partial_x q = A_1^2 + y A_3^2, \qquad
  \partial_y q = A_2^2 + x A_3^2, \qquad
  Dq = A_3^2 .
\end{equation}
So $Dq$ is always a perfect square, and $Dq = 0$ exactly when $q$ has
no monomial with both exponents odd --- equivalently, when
$q \in \F_2[x^2, y] + \F_2[x, y^2]$, i.e.\ when $q$ partially factors
through the squaring map.  The polynomial $A_3$ measures the failure
of this degeneration, and the whole proof is organized around it.

The degenerate case $Dq = 0$ is settled by a conserved quantity:

\begin{lemma}[chain rule for $D$]\label{lem:cocycle}
In characteristic $2$, for all $q, \alpha, \beta \in \F_2[x,y]$:
\[
  D\bigl(q(\alpha,\beta)\bigr)
  = (Dq)(\alpha,\beta)\,
      \bigl(\partial_x\alpha\,\partial_y\beta
            + \partial_y\alpha\,\partial_x\beta\bigr)
  + (\partial_x q)(\alpha,\beta)\, D\alpha
  + (\partial_y q)(\alpha,\beta)\, D\beta .
\]
\end{lemma}

\begin{proof}[Proof sketch]
Expand $\partial_x\partial_y\bigl(q(\alpha,\beta)\bigr)$ by the
ordinary chain and product rules.  The terms that would involve
$\partial_x^2 q$ or $\partial_y^2 q$ vanish by fact 2, and the
survivors assemble as displayed.  (In characteristic $0$ the same
computation produces additional second-derivative terms; their absence
here is the point.)
\end{proof}

\begin{proposition}\label{prop:Dzero}
If $Dq = 0$ then $xy \notin \Clo(q)$.
\end{proposition}

\begin{proof}
The set $\{g : Dg = 0\}$ contains the variables and constants, and by
Lemma~\ref{lem:cocycle} (with the first term dead) it is closed under
substitution into $q$.  Hence it contains all of $\Clo(q)$.  But
$D(xy) = 1$.
\end{proof}

From now on suppose $Dq \neq 0$, i.e.\ $A_3 \neq 0$; we must show
$x + y \notin \Clo(q)$.

\subsection{The rigidity property $(\dagger)$ and the reduction to
it}\label{sec:tameness}

Write $\F_2[x{+}y]$ for the subring of polynomials in $x+y$ (the
``diagonal'' polynomials).  The proof of Theorem~\ref{thm:F2} in the
case $Dq \neq 0$ runs through the following property of $q$:
\begin{quote}
  $(\dagger)$\quad for all $\alpha, \beta \in \F_2[x,y]$: whenever
  $q(\alpha,\beta)$ is a \emph{nonconstant} element of
  $\F_2[x{+}y]$, both $\alpha, \beta \in \F_2[x{+}y]$.
\end{quote}
In words: the only way a substitution into $q$ can land on a
nonconstant diagonal polynomial is for both inputs to be diagonal
already.  The property is designed to feed the following induction.

\begin{proposition}[from $(\dagger)$ to non-generation]\label{prop:master}
If $q$ satisfies $(\dagger)$ then $x + y \notin \Clo(q)$.
\end{proposition}

\begin{proof}
We claim every $g \in \Clo(q)$ satisfies
\[
  P(g): \qquad g \notin \F_2[x{+}y]
  \quad\text{or}\quad g \text{ is constant.}
\]
The variables satisfy $P$: the substitution $x \leftrightarrow y$
fixes every element of $\F_2[x{+}y]$ but moves $x$, so
$x \notin \F_2[x{+}y]$, likewise $y$.  Constants satisfy $P$.
Inductive step: let $g = q(\alpha,\beta)$ with $P(\alpha)$,
$P(\beta)$, and suppose $P(g)$ fails: $g \in \F_2[x{+}y]$,
nonconstant.  Property $(\dagger)$ forces
$\alpha, \beta \in \F_2[x{+}y]$; then
$P(\alpha), P(\beta)$ force $\alpha, \beta$ constant; but then $g$ is
constant, a contradiction.  So $P$ holds on the clone, and $x+y$ ---
which is in $\F_2[x{+}y]$ and nonconstant --- is not in it.
\end{proof}

So Theorem~\ref{thm:F2} reduces to:

\begin{theorem}\label{thm:tame}
If $Dq \neq 0$ then $q$ satisfies $(\dagger)$.
\end{theorem}

\subsection{From polynomials to curves}\label{sec:bridge}

Property $(\dagger)$ becomes a geometric statement after restriction
to a generic member of the pencil of lines $\{x + y = \text{const}\}$.  Fix once
and for all an algebraically closed field $K \supseteq \F_2$ of
characteristic $2$ containing an element $\sigma$ transcendental over
$\F_2$; concretely, $K = \overline{\F_2(\sigma)}$, an algebraic
closure of the field of rational functions in one indeterminate.  We
say $c \in K$ is \emph{transcendental} if it satisfies no nonzero
polynomial relation over $\F_2$ --- the algebraic stand-in for ``$c$
is a generic value''.

\begin{definition}\label{def:witness}
A \emph{witness} is a pair $\alpha, \beta \in K[t]$, not both
constant, such that $q(\alpha, \beta) = c$ is a constant, with $c$
transcendental.
\end{definition}

Geometrically: a polynomial-parametrized curve
$t \mapsto (\alpha(t), \beta(t))$ in the plane over $K$, lying
entirely inside the single level set $\{q = c\}$ of a generic level
$c$.

\begin{proposition}[from curves to $(\dagger)$]\label{prop:bridge}
If no witness exists, then $q$ satisfies $(\dagger)$.
\end{proposition}

\begin{proof}
Let $\alpha, \beta \in \F_2[x,y]$ with $q(\alpha,\beta) = f(x+y)$,
$f \in \F_2[u]$ nonconstant.  Substitute the parametrized generic line
$(x, y) = (t,\, t + \sigma)$, under which $x + y \mapsto \sigma$.  The
left side becomes
$q\bigl(\alpha(t, t{+}\sigma),\, \beta(t, t{+}\sigma)\bigr)$ and the
right side the \emph{constant} $f(\sigma)$, which is transcendental
(a nonzero $\F_2$-relation on $f(\sigma)$ composes with $f$ to a
nonzero relation on $\sigma$).  Since no witness exists, both
$\alpha(t, t{+}\sigma)$ and $\beta(t, t{+}\sigma)$ are constant in
$t$.  It remains to see that a polynomial $g \in \F_2[x,y]$ that is
constant along the generic line lies in $\F_2[x{+}y]$: expand
$g(x,\, x + u) = \sum_i g_i(u)\, x^i$ with $g_i \in \F_2[u]$;
constancy in $t$ of $g(t, t+\sigma)$ says $g_i(\sigma) = 0$ for all
$i \geq 1$, hence $g_i = 0$ ($\sigma$ transcendental), so
$g(x, x+u) = g_0(u)$, i.e.\ $g = g_0(x + y)$ after substituting
$u = y - x$ (which in characteristic 2 is $y + x$).
\end{proof}

Theorem~\ref{thm:tame} has thus become a concrete statement we can
attack with one-variable calculus:

\begin{theorem}[no witnesses]\label{thm:descent}
If $Dq \neq 0$, no witness exists.
\end{theorem}

\section{The Frobenius descent}\label{sec:descent}

This section proves Theorem~\ref{thm:descent}.  We first develop the
mechanism in the simplest nontrivial situation, then do the general
case.

Throughout, $(\alpha, \beta)$ is a putative witness:
$\alpha, \beta \in K[t]$ not both constant, $q(\alpha,\beta) = c$,
$c$ transcendental.  A prime $'$ denotes $d/dt$.  Degrees are degrees
in $t$, under the convention of Section~\ref{sec:prelim} (constants,
including the zero polynomial, have degree $0$).  Recall the two
calculus facts, now over $K[t]$: squares have zero derivative; and
conversely, since $K$ is algebraically closed (hence \emph{perfect}:
every element has a square root), a polynomial in $K[t]$ with zero
derivative is a square --- it is a polynomial in $t^2$, and its
coefficients have square roots.

\subsection{Warm-up: the case $A_3 = 1$}\label{sec:warmup}

Suppose first the parity decomposition \eqref{eq:parity} has
$A_3 = 1$:
\[
  q = A_0^2 + x A_1^2 + y A_2^2 + xy,
  \qquad
  \partial_x q = A_1^2 + y, \quad \partial_y q = A_2^2 + x .
\]
(The smallest example is $q = xy$ itself; the case already contains
all polynomials whose unique both-exponents-odd monomial is $xy$.)

Differentiate the witness equation $q(\alpha,\beta) = c$ in $t$.  The
chain rule gives
\begin{equation}\label{eq:chain-warm}
  A\,\alpha' + B\,\beta' = 0,
  \qquad
  A := (\partial_x q)(\alpha,\beta), \quad
  B := (\partial_y q)(\alpha,\beta).
\end{equation}
Next differentiate $A$ and $B$ themselves.  Using
$\partial_x(\partial_x q) = 0$ and
$\partial_y(\partial_x q) = Dq = 1$ (and symmetrically):
\begin{equation}\label{eq:engine-warm}
  A' = \beta', \qquad B' = \alpha' .
\end{equation}
Combining \eqref{eq:chain-warm} and \eqref{eq:engine-warm}:
\[
  (AB)' = A'B + AB' = \beta' B + A \alpha' = 0,
\]
so $AB$ is a \emph{square} in $K[t]$: along a
witness, the product of the two partial derivatives is forced to be a
perfect square.

Now the two cases.  (For this warm-up we take for granted that neither $A$
nor $B$ vanishes identically; this degeneracy, like several others in
the general case, is dispatched by Lemma~\ref{lem:kill} below,
and we do not belabor it here.)  Suppose
first $A$ and $B$ have a common root $\theta \in K$.  Then the point
$P := (\alpha(\theta), \beta(\theta))$ satisfies
$\partial_x q (P) = \partial_y q(P) = 0$: it is a \emph{critical
point} of $q$.  In the warm-up case the critical equations read
$A_1^2(P) = \beta(\theta)$ and $A_2^2(P) = \alpha(\theta)$ --- two
polynomial equations over $\F_2$ pinning $P$ down; one shows (this is
easy here, and is Lemmas~\ref{lem:algpoint} and~\ref{lem:keystone} in
general) that $P$ must have coordinates \emph{algebraic} over $\F_2$.
But then $c = q(P)$ is algebraic --- contradicting transcendence of
the level.  A generic level set passes through no critical point.

Otherwise $A$ and $B$ are coprime in $K[t]$.  A product of two coprime
polynomials is a square only if each factor is (up to unit, and units
of $K[t]$ are constants of the algebraically closed $K$, hence
squares).  So $A$ and $B$ are squares; squares have zero derivative;
and \eqref{eq:engine-warm} turns this into
\[
  \alpha' = \beta' = 0 .
\]
By perfectness, $\alpha = \alpha_1^2$ and $\beta = \beta_1^2$ for some
$\alpha_1, \beta_1 \in K[t]$ of \emph{half} the degree.  Here the
$\F_2$-coefficients enter: squaring fixes $0$ and $1$, so
squaring a polynomial commutes with substituting into $q$,
\[
  q(\alpha_1, \beta_1)^2
  \;=\; q(\alpha_1^2,\, \beta_1^2)
  \;=\; q(\alpha, \beta) \;=\; c .
\]
Writing $r := \sqrt{c} \in K$ (transcendental, since $r$ algebraic
would make $c = r^2$ algebraic), we conclude
$q(\alpha_1, \beta_1) = r$: \emph{the pair $(\alpha_1, \beta_1)$ is
again a witness, for the same $q$, at half the total degree}.

Iterating, the total degree $\deg\alpha + \deg\beta$ halves at every
round in the second branch, and cannot do so forever; the first branch
must eventually occur, giving the contradiction.  This proves
Theorem~\ref{thm:descent} in the warm-up case.  The induction
variable, to repeat the point from the introduction, is a property of
the \emph{witness} --- how many more times it can factor through
squaring, its Frobenius depth --- while $q$ stands still.

\subsection{What changes in general}

For general $q$ with $Dq = A_3^2 \neq 0$, two new things happen.

First, the derivative identities \eqref{eq:engine-warm} acquire the factor
$\Delta := (Dq)(\alpha,\beta) = A_3(\alpha,\beta)^2$:
\[
  A' = \Delta\,\beta', \qquad B' = \Delta\,\alpha',
\]
which is harmless \emph{provided} $\Delta \neq 0$ along the witness
--- a new way the argument can degenerate.

Second and more seriously, the first of the two cases needs critical
points of $q$ to be algebraic, and this can simply be false: the
critical equations $\partial_x q = \partial_y q = 0$ may have
\emph{curves} of solutions.  Indeed if $w := \gcd(A_1, A_2, A_3)$ is
nonconstant then by \eqref{eq:partials} the whole curve $\{w = 0\}$
consists of critical points, and a witness might pass through such
points.  The remedy is to peel $w$ off before running the
argument.  Write $A_i = w B_i$, so that the $B_i$ have no common factor,
and
\begin{equation}\label{eq:peel}
  q = A_0^2 + w^2\, h,
  \qquad
  h := x B_1^2 + y B_2^2 + xy B_3^2 ,
\end{equation}
with $B_3 \neq 0$; we call the gcd extraction \eqref{eq:peel} the
\emph{peel}.  We will run the warm-up argument \emph{on $h$
instead of $q$}: the function
$H := h(\alpha,\beta)$ is no longer constant along the witness, but
its derivative still vanishes, which is what makes the peel legal.
From \eqref{eq:peel} and
$q(\alpha,\beta) = c = r^2$:
\begin{equation}\label{eq:S}
  W^2 H = S^2,
  \qquad
  W := w(\alpha,\beta), \quad S := r + A_0(\alpha,\beta),
\end{equation}
using $(u+v)^2 = u^2 + v^2$.  If $W \neq 0$, differentiating
\eqref{eq:S} (both $W^2$ and $S^2$ are squares) gives $W^2 H' = 0$,
so $H' = 0$.  The
derivative argument only ever used the vanishing of the level's
derivative, so it runs on $h$ unchanged.

What the peel buys is the algebraicity of critical points: for $h$,
unlike $q$, the critical locus contains no curve.  That is
Lemma~\ref{lem:keystone} below.  Finally, the various nondegeneracy conditions accumulated
along the way ($W \neq 0$, $\Delta \neq 0$, $A \neq 0$, $B \neq 0$)
are all of the same shape --- some fixed nonzero polynomial with
$\F_2$-coefficients vanishing identically along the witness --- and a
single lemma disposes of them all: no such relation can hold, because
it would tie the generic level to algebraic data.

We now state the three lemmas and assemble the proof.

\subsection{Three lemmas}

\begin{lemma}[no algebraic relation along a witness]\label{lem:kill}
Let $(\alpha,\beta)$ be a witness with level $c$, and
$f \in \F_2[x,y]$, $f \neq 0$.  Then $f(\alpha,\beta) \neq 0$ in
$K[t]$.
\end{lemma}

\begin{proof}[Proof sketch]
Suppose $f(\alpha,\beta) = 0$; after swapping variables if needed,
assume $\alpha$ nonconstant.  Adjoin a fresh variable $z$ standing for
the level, and consider $f$ and $g := q + z$ as polynomials in $y$
over the ring $\F_2[x,z]$.  Since $g$ is monic and linear in $z$, it
is irreducible and cannot divide $f$ (which is $z$-free), so $f$ and
$g$ are coprime; a B\'ezout identity over the fraction field, with
denominators cleared, produces
\[
  f\,U + (q + z)\,V = E(x, z) \neq 0
  \qquad \text{in } \F_2[x,z][y].
\]
Now substitute $y \mapsto \beta$, $x \mapsto \alpha$, $z \mapsto c$:
the first summand dies by hypothesis, the second because
$q(\alpha,\beta) + c = c + c = 0$ in characteristic $2$.  So
$E(\alpha, c) = 0$.  Expanding $E = \sum_i \rho_i(z)\,x^i$ and using
that $1, \alpha, \alpha^2, \dots$ have strictly increasing degrees in
$t$ (here nonconstancy of $\alpha$ enters), each coefficient
$\rho_i(c)$ must vanish; some $\rho_i \neq 0$ then exhibits $c$ as
algebraic over $\F_2$ --- contradiction.
\end{proof}

\begin{lemma}[algebraic points]\label{lem:algpoint}
If $u, v \in \F_2[x,y]$ have no common nonunit factor and
$(a,b) \in K^2$ satisfies $u(a,b) = v(a,b) = 0$, then $a, b$ are
algebraic over $\F_2$.
\end{lemma}

\begin{proof}[Proof sketch]
Viewing $u, v$ in $\F_2(x)[y]$: a common divisor there would descend
(Gauss's lemma; see e.g.\ \cite[Ch.~IV]{Lang}) to a common factor in
$\F_2[x,y]$, so they are
coprime, and a cleared B\'ezout identity gives
$su + tv = r(x) \neq 0$; evaluate at $(a,b)$ to get $r(a) = 0$.
Symmetrically for $b$.
\end{proof}

\begin{lemma}[coprimality of the peeled
partials]\label{lem:keystone}
Let $B_1, B_2, B_3 \in \F_2[x,y]$ have no common nonunit factor, with
$B_3 \neq 0$, and set
$h_x := B_1^2 + y B_3^2$, $h_y := B_2^2 + x B_3^2$ (the partials of
$h$).  Then $h_x$ and $h_y$ have no common nonunit factor.
\end{lemma}

\begin{proof}
Both are nonzero: applying $\partial_y$ to a relation $h_x = 0$ would
give $B_3^2 = 0$.  Suppose a prime $p$ divides both.

If $p \mid B_3$: subtracting the $B_3^2$-terms, $p$ divides $B_1^2$
and $B_2^2$, hence (primality) all three $B_i$ --- excluded.

If $p \nmid B_3$: write $h_x = p u$ and $h_y = p v$.  Differentiate
$h_x = B_1^2 + y B_3^2$, remembering squares are derivative-free:
\[
  0 = \partial_x h_x = p_x u + p\, u_x
  \quad\Longrightarrow\quad p \mid p_x u ;
\]
\[
  B_3^2 = \partial_y h_x = p_y u + p\, u_y
  \quad\Longrightarrow\quad p \nmid u
\]
(else $p \mid B_3^2$, so $p \mid B_3$);
\[
  B_3^2 = \partial_x h_y = p_x v + p\, v_x
  \quad\Longrightarrow\quad p \nmid p_x
\]
(same reason).  So the prime $p$ divides the product $p_x u$ but
neither factor --- contradiction.
\end{proof}

\begin{remark}
Lemma~\ref{lem:keystone} has geometric content: a curve of common
zeros of $h_x, h_y$ would have a function field in which both $x$ and
$y$ become squares (solve the critical equations for $y$ and $x$),
making that function field \emph{perfect}; but the function field of a
curve in characteristic $2$ is never perfect.  The displayed proof is
this argument in infinitesimal form --- ``$p \nmid p_x$'' is the
derivation that survives --- and entirely avoids the geometry.
\end{remark}

\subsection{Proof of Theorem~\ref{thm:descent}}

Among all witnesses, take one minimizing
$n := \deg\alpha + \deg\beta$ ($\geq 1$, as not both are constant).
Decompose $q$ as in \eqref{eq:peel}; fix $r = \sqrt{c}$,
transcendental.

By Lemma~\ref{lem:kill} applied to each of the
nonzero polynomials $w$, $h_x$, $h_y$, $B_3$ --- nonzero by
construction and by the argument opening Lemma~\ref{lem:keystone} ---
none of
\[
  W := w(\alpha,\beta), \quad
  A := h_x(\alpha,\beta), \quad
  B := h_y(\alpha,\beta), \quad
  \Delta := B_3(\alpha,\beta)^2
\]
vanishes in $K[t]$.

\emph{Peel.}  As in \eqref{eq:S}, $W^2 H = S^2$ with
$H = h(\alpha,\beta)$, $S = r + A_0(\alpha,\beta)$, and
differentiating, $H' = 0$.

\emph{Derivative identities.}  The chain rule for $h$, plus
$\partial_x h_x = 0$,
$\partial_y h_x = \partial_x h_y = Dh = B_3^2$,
$\partial_y h_y = 0$, gives
\[
  A\alpha' + B\beta' = H' = 0,
  \qquad A' = \Delta\beta', \qquad B' = \Delta\alpha' ,
\]
whence $(AB)' = \Delta(A\alpha' + B\beta') = 0$ and $AB$ is a square.

\emph{Case 1: $A$ and $B$ share a root $\theta$.}  Then
$P := (\alpha(\theta), \beta(\theta))$ is a common zero of
$h_x, h_y$, which by Lemmas~\ref{lem:algpoint}
and~\ref{lem:keystone} has algebraic coordinates.  Evaluate the peel
identity \eqref{eq:S} at $\theta$:
\[
  S(\theta)^2 = w(P)^2\, h(P)
\]
is algebraic, hence so is $S(\theta)$, hence so is
$r = S(\theta) + A_0(P)$.  This contradicts transcendence of $r$.

\emph{Case 2: $A, B$ coprime.}  Coprime factors of a square
are squares, so $A' = B' = 0$; since $\Delta \neq 0$, the derivative
identities force $\alpha' = \beta' = 0$, so $\alpha = \alpha_1^2$,
$\beta = \beta_1^2$, and as in the warm-up
$q(\alpha_1, \beta_1) = r$.  Then $(\alpha_1, \beta_1)$ is a witness
with $\deg\alpha_1 + \deg\beta_1 = n/2 < n$, contradicting
minimality.  (If $n = 1$ this branch cannot occur at all, as halving
is impossible; so the descent is well-founded.)  \qed

\medskip

Assembling: Theorem~\ref{thm:descent} $\Rightarrow$
(Proposition~\ref{prop:bridge}) property $(\dagger)$ $\Rightarrow$
(Proposition~\ref{prop:master})
$x+y \notin \Clo(q)$ when $Dq \neq 0$; together with Proposition~\ref{prop:Dzero} for
$Dq = 0$, this proves Theorem~\ref{thm:F2}, and with
Lemma~\ref{lem:transfer}, Theorem~\ref{thm:Z}.  \qed

\section{From $\F_2$ to the dichotomy}\label{sec:tothedichotomy}

Theorem~\ref{thm:F2} is the kernel of the negative half; the rest of
Theorem~\ref{thm:dichotomy} follows from it by three short
reductions, which we now record.

\medskip
\noindent\emph{Step 1: perfect fields, via a coefficient twist.}

\begin{proposition}\label{prop:perfect}
Let $k$ be a perfect field of characteristic $2$.  For every
$q \in k[x,y]$, at least one of $x+y$, $xy$ lies outside $\Clo(q)$.
\end{proposition}

\begin{proof}
The descent generalizes to any perfect coefficient field $k$ of
characteristic $2$ with one modification.  The coefficients of $q$
are no longer fixed by squaring, so in the halving branch the
identity $q(\alpha_1^2, \beta_1^2) = \tilde q(\alpha_1, \beta_1)^2$
holds not for $q$ itself but for the \emph{coefficient-twisted}
polynomial $\tilde q = q^{(1/2)}$ (apply the inverse Frobenius to
each coefficient --- this is where perfectness enters).  The halving
branch therefore produces a half-degree witness for $q^{(1/2)}$
rather than for $q$; since the induction is on the witness degree
with $q$ universally quantified inside, the twist is absorbed at no
cost, and no periodicity argument is needed.  The remaining
ingredients --- the parity decomposition and
Propositions~\ref{prop:master} and~\ref{prop:bridge} --- are
unaffected, and the disjunctive conclusion follows
as over $\F_2$.
\end{proof}

\medskip
\noindent\emph{Step 2: arbitrary fields of characteristic $2$, via
the perfection.}

\begin{proof}[Proof of Theorem~\ref{thm:char2}]
An arbitrary field $k$ of characteristic $2$ embeds into its
perfection $k^{1/2^\infty}$, which is perfect of characteristic $2$.
The coefficient embedding is a ring homomorphism, so by the transfer
lemma (Lemma~\ref{lem:transfer}) a derivation of $x+y$ (or $xy$) over
$k$ maps to one over the perfection; if both $x+y$ and $xy$ lay in
$\Clo(q)$ over $k$, both would lie in the clone over the perfection,
contradicting Proposition~\ref{prop:perfect}.
\end{proof}

\medskip
\noindent\emph{Step 3: rings in which $2$ is not a unit, via residue
fields.}

\begin{proof}[Proof of Theorem~\ref{thm:dichotomy}]
If $2$ is a unit of $R$, the operation $x^2 - y$ generates $R[x,y]$
by Theorem~\ref{thm:converse}.  If $2$ is not a unit, it lies in a
maximal ideal $\mathfrak m$, and a generator over $R$ would
push forward along $R \to R/\mathfrak m$ (Lemma~\ref{lem:transfer}
again) to a generator over the residue field, which has
characteristic $2$ --- contradicting Theorem~\ref{thm:char2}.
\end{proof}

\section{Relation to classical functional completeness}
\label{sec:context}

For operations-as-functions on a finite set, one-element functional
completeness
is classical territory.  Post's classification of the Boolean clones
locates the Sheffer stroke \cite{Sheffer,Post}; Webb exhibited a
single functionally complete operation on every finite domain
\cite{Webb}; and
there are effective functional completeness criteria, from S{\l}upecki through
Rousseau to Rosenberg's classification of the maximal function clones
\cite{Slupecki,Rousseau,Rosenberg}.

Example~\ref{ex:nand} (Section~\ref{sec:precision}) showed that over
finite fields the formal statements of this paper and the classical
functional ones can diverge.  In the positive direction, however,
the generation theorem does project onto a classical statement,
which we record for universal algebraists.
Recall that an algebra is \emph{functionally complete} if every
finitary operation on its underlying set is a polynomial operation
(= a term operation after all constants are adjoined).

\begin{corollary}\label{cor:funcomplete}
Let $F$ be a finite field of odd order.  Then the algebra
$\langle F;\; x^2 - y \rangle$ with all constants adjoined is
functionally complete: every function $F^n \to F$, for every $n$, is
a composite of $x^2 - y$, constants, and variables.
\end{corollary}

\begin{proof}
By Theorem~\ref{thm:converse} (proved in Section~\ref{sec:converse};
$2$ is invertible in $F$), $x + y$ and $xy$ are composites of
$x^2 - y$ and constants, hence so is every polynomial in any number
of variables; and by Lagrange interpolation every function
$F^n \to F$ is induced by a polynomial.
\end{proof}

\noindent
Functional completeness criteria for finite algebras are classical
territory, from Foster's primal algebras and the Foster--Pixley
theory of semi-primal algebras \cite{Foster,FosterPixley} to Werner's
discriminator varieties \cite{Werner}; Corollary~\ref{cor:funcomplete}
places $\langle F;\, x^2 - y\rangle$ (with constants) inside that
landscape, whereas Theorem~\ref{thm:F2}, by Example~\ref{ex:nand},
lives strictly outside it.

\medskip

The composition structure of polynomial rings has its own literature
--- Menger's tri-operational algebras, Adler's composition rings, and
near-ring theory for univariate composition \cite{Menger,Adler,Pilz},
as well as Ritt's theory of decompositions of a single univariate
polynomial \cite{Ritt}.  These study the composition structure of the
full polynomial ring rather than generation from a single binary
operation, and we have not found the present question addressed
there.  We would be grateful for pointers to antecedents; the
question is natural enough that we hesitate to claim novelty for
anything but the answers.

\section{A note on formal verification}\label{sec:lean}

The proofs above, especially the descent, involve a large amount of
routine characteristic-$2$ manipulation, and we wanted an independent
check.  All results stated here --- Theorems~\ref{thm:Z}, \ref{thm:dichotomy}, \ref{thm:converse},
\ref{thm:F2} and the supporting lemmas --- have
been verified in the Lean~4 proof assistant against the Mathlib
library, with no unproved assumptions beyond Lean's standard
foundational axioms.  The formal statements involve only the
fifteen-line definition of $\Clo$ reproduced in
Appendix~\ref{app:lean}.  In addition, the headline theorems are
certified by an adversarial judge (\texttt{leanprover/comparator}):
their statements live in a standalone challenge file of roughly
ninety lines that does not import the development, and the judge
replays all proofs through two independent proof kernels against
those exact statements, so trusting the results requires reading only
the challenge file and the verdict.  Two of the proofs presented above are
actually simplifications found during the verification process: the
derivative proof of Lemma~\ref{lem:keystone}, which
replaced a longer argument via perfect function fields, and the
B\'ezout-certificate form of Lemma~\ref{lem:kill}.  We will not say more
about the formalization here; the development --- including a
blueprint mapping each informal statement of this paper to its formal
declaration --- is available from the author and will be made public
with the paper.

\section{Open questions}\label{sec:open}

\subsection*{Degrees and characterization}

\begin{question}
Which degrees admit single generators over $\Q$?  Degree $2$ does
(Theorem~\ref{thm:converse}).  For degree $3$ a natural candidate is
$x^3 - y$: translations by arbitrary rationals follow from the
classical fact that every rational number is a sum of three rational
cubes (with explicit formulas), and a cubic analogue of the linear
span step appears to reach $xy$; we have not finished the
verification.  The answer is characteristic-sensitive: over $\F_3$
the map $u \mapsto u^3$ is additive (Frobenius), and one checks that
$\Clo(x^3 - y)$ consists only of polynomials of the form
\emph{additive plus constant} --- $\F_3$-linear combinations of the
monomials $x^{3^i}, y^{3^j}$, plus a constant (the constants enter
because they are generators) --- a class closed under substitution
into $x^3 - y$ and excluding $xy$; so $x^3 - y$ is certainly not a
generator there.
\end{question}

\begin{question}\label{q:characterize}
Away from characteristic $2$, \emph{which} operations are single
generators?
Generation is not simply ``nonlinear and depending on both
variables'': over $\Q$,
\begin{itemize}
\item $xy + 1$ is not a generator, although it is nonlinear, uses both
  variables, and $\Clo(xy+1)$ contains $xy$.  The obstruction is a
  characteristic-zero instance of the rigidity property $(\dagger)$:
  if $\alpha\beta + 1 = F(x{+}y)$ is nonconstant then
  $\alpha\beta = h(x{+}y)$ splits over $\overline{\Q}$ into linear
  factors $(x{+}y-r_i)$, so by unique factorization every divisor ---
  in particular $\alpha$ and $\beta$ --- lies in $\Q[x{+}y]$;
  Proposition~\ref{prop:master} (whose proof is characteristic-free)
  then excludes $x+y$.
\item $x^2 + y$ is not a generator, although $x^2 - y$ is: every element
  of $\Clo(x^2+y)$ has the form $a + \sum_i s_i^2$ with
  $a \in \{x, y\} \cup \Q$ (the shape is preserved by
  $(f,g) \mapsto f^2 + g$), and neither $x+y$ nor $xy$ is of this
  form, since sums of squares of real polynomials take no negative
  values.  A \emph{positivity} obstruction from the real place,
  invisible over finite fields.
\end{itemize}
Together with the affine obstruction (linear $p$) and the monomial
obstruction ($\Clo(xy)$ consists of monomials), this suggests asking
whether generation is governed by a finite explicit list of
``conserved classes'' --- proper subsets of $R[x,y]$ containing the
generators and closed under substitution into $p$.  We do not know;
we also cannot presently decide simple boundary cases such as
$x^2 - 2y$, where the increment $\gamma \mapsto 4\gamma +
(\alpha^2 - 2\alpha'^2)$ rescales instead of translating and a
$2$-adic conserved class is conceivable.  A positive structure
theorem here would likely also bear on
Question~\ref{q:decidable}.
\end{question}

\subsection*{Decidability}

\begin{question}\label{q:decidable}
Is membership in $\Clo(p)$ decidable?  Is generation?
\end{question}

Half of the problem is easy: membership is semidecidable, since one
can enumerate all formulas in $p$ and expand each.  It is the
complementary half that we want to dwell on, because the difficulty
is structural and shaped our whole approach.  The obvious idea ---
search all formulas up to some size or degree bound, and conclude
non-membership if the target never appears --- proves nothing, no
matter the bound, because substitution is not monotone in any evident
complexity measure.  The culprit is \emph{cancellation}.  Already for
$q = x^2 + y$ over $\F_2$ one has, for \emph{all} $f, g$,
\[
  q\bigl(f,\; q(f, g)\bigr) \;=\; f^2 + (f^2 + g) \;=\; g :
\]
the intermediate polynomial $f$, of arbitrarily large degree, is
annihilated without trace.  So every element of the clone admits
derivations passing through intermediates of unbounded degree, and a
small target may, for all a bounded search can tell, first become
reachable only after a long excursion through enormous polynomials
that cancel at the last step.  The degree of what you are building
says nothing about the degree of what you must build it from.

This is why every non-membership argument in this paper avoids search
entirely and instead exhibits a \emph{cancellation-immune invariant}:
a property that is checked on the final polynomial alone and is
preserved by every substitution into $q$, whatever the intermediate
terms are --- membership in $\ker D$ in
Proposition~\ref{prop:Dzero}, the invariant $P$ in
Proposition~\ref{prop:master}.  A decision procedure would seem to
require either a computable bound on derivation size, which
cancellation undermines, or a supply of such invariants
rich enough to certify every non-membership; we possess neither,
and we suspect the question is hard.  For comparison: in the cousin
problem of membership in the \emph{ideal} generated by given
polynomials, the analogous degree question is settled --- doubly
exponential degree bounds make ideal membership decidable, indeed
EXPSPACE-complete by Mayr--Meyer \cite{MayrMeyer} --- and it is
exactly such a bound that cancellation withholds here.  As a
concrete instance of how
little we know: for $q$ with $Dq = 0$, Theorem~\ref{thm:F2} settles
$xy \notin \Clo(q)$, but we know no algorithm deciding, for such $q$,
whether $x + y \in \Clo(q)$.

\subsection*{Universal algebra}

\begin{question}\label{q:ua}
The universal-algebra side of the picture is wide open beyond
Corollary~\ref{cor:funcomplete}, which adjoins all constants.
Without constants: is $\langle F;\, x^2 - y \rangle$ \emph{primal}
for every finite field $F$ of odd order --- is every finitary
operation on $F$ a term operation of $x^2 - y$ alone?  (Rousseau's
criterion \cite{Rousseau} reduces this to a concrete finite check for
any one $F$; we have not pursued it.)  More broadly, we have
attempted no Mal'cev-condition analysis of the algebras
$\langle R;\, p \rangle$: which congruence conditions
(permutability, distributivity, \dots) the varieties they generate
satisfy, and whether these interact with $p$ generating $R[x,y]$, are
natural questions for universal algebraists that our methods do not
touch.
\end{question}

\subsection*{Rational functions}

\begin{question}
The same questions can be posed for the field of rational functions:
preliminary calculations suggest $1/x - y$ is a single generator of
$\Q(x,y)$ in a suitable sense.  Substitution of rational functions is
a partial operation (denominators may vanish identically), so even
stating this precisely requires a partial-clone formalism.
\end{question}

\medskip
\noindent\textbf{Draft note.}  All references below have been checked
against the published record, with three exceptions, each marked in
the source: the issue and page numbers of the Menger item could not
be confirmed against the original (the series and year are
confirmed); the journal volume and pages of the Goldstern--Pinsker
item were checked only against secondary sources (the arXiv version
is confirmed); and the published version of the Aichinger--Kreinecker
item could not be fully reconciled with the arXiv version, which is
the one cited.  Further pointers, especially to antecedents of the
main question, remain welcome.

\begin{thebibliography}{99}

\bibitem{Adler} I.~Adler, \emph{Composition rings}, Duke Math.\ J.\
29 (1962), 607--623.

\bibitem{Aichinger2013} E.~Aichinger, \emph{Generating polynomials
using function composition}, Quaest.\ Math.\ 36 (2013), no.~1,
39--46.

% UNVERIFIED: the published version appears as S. Kreinecker,
% "Generating integer polynomials from X^2 and X^3 using function
% composition: a study of subnearrings of (Z[X], +, o)", Quaest.
% Math. 44 (2021), no. 5, 693--715, with an author list differing
% from the arXiv version; this was checked only against search
% snippets, so we cite the (confirmed) arXiv version.
\bibitem{AichingerKreinecker} E.~Aichinger and S.~Kreinecker,
\emph{Subnearrings of $(\Z[x], +, \circ)$}, preprint,
arXiv:1709.01345 (2017).

\bibitem{Foster} A.~L.~Foster, \emph{Generalized ``Boolean'' theory
of universal algebras.  Part II.  Identities and subdirect sums of
functionally complete algebras}, Math.\ Z.\ 59 (1953), 191--199.

\bibitem{FosterPixley} A.~L.~Foster and A.~F.~Pixley,
\emph{Semi-categorical algebras.  I.  Semi-primal algebras}, Math.\
Z.\ 83 (1964), 147--169.

% UNVERIFIED: journal volume and page numbers checked only against
% secondary sources; the arXiv version is confirmed.
\bibitem{GoldsternPinsker} M.~Goldstern and M.~Pinsker, \emph{A
survey of clones on infinite sets}, Algebra Universalis 59 (2008),
365--403; arXiv:math/0701030.

\bibitem{HamkinsMO} J.~D.~Hamkins, \emph{Can we unify addition and
multiplication into one binary operation?  To what extent can we
find universal binary operations?}, MathOverflow question 57465
(2011), \url{https://mathoverflow.net/q/57465}.

\bibitem{KaarliPixley} K.~Kaarli and A.~F.~Pixley, \emph{Polynomial
Completeness in Algebraic Systems}, Chapman \& Hall/CRC, Boca Raton,
2001.

\bibitem{Lang} S.~Lang, \emph{Algebra}, rev.\ 3rd ed., Graduate Texts
in Mathematics 211, Springer, New York, 2002.

\bibitem{Lawvere} F.~W.~Lawvere, \emph{Functorial semantics of
algebraic theories}, Proc.\ Nat.\ Acad.\ Sci.\ USA 50 (1963),
869--872.

\bibitem{MayrMeyer} E.~W.~Mayr and A.~R.~Meyer, \emph{The complexity
of the word problems for commutative semigroups and polynomial
ideals}, Adv.\ in Math.\ 46 (1982), 305--329.

% UNVERIFIED: issue and page numbers of the Menger entry could not be
% confirmed against the original (series and year confirmed).
\bibitem{Menger} K.~Menger, \emph{Tri-operational algebra}, Reports
of a Mathematical Colloquium (2) 5/6 (1944), 3--10.

\bibitem{Odrzywolek} A.~Odrzywo{\l}ek, \emph{All elementary
functions from a single binary operator}, preprint, arXiv:2603.21852
(2026).

\bibitem{Pilz} G.~Pilz, \emph{Near-Rings: The Theory and its
Applications}, rev.\ ed., North-Holland Mathematics Studies 23,
North-Holland, Amsterdam, 1983.

\bibitem{Post} E.~L.~Post, \emph{The Two-Valued Iterative Systems of
Mathematical Logic}, Annals of Mathematics Studies 5, Princeton
University Press, Princeton, 1941.

\bibitem{Ritt} J.~F.~Ritt, \emph{Prime and composite polynomials},
Trans.\ Amer.\ Math.\ Soc.\ 23 (1922), 51--66.

\bibitem{Rosenberg} I.~G.~Rosenberg, \emph{\"Uber die funktionale
Vollst\"andigkeit in den mehrwertigen Logiken.  Struktur der
Funktionen von mehreren Ver\"anderlichen auf endlichen Mengen},
Rozpravy \v{C}eskoslovensk\'e Akad.\ V\v{e}d \v{R}ada Mat.\
P\v{r}\'irod.\ V\v{e}d 80 (1970), no.~4, 3--93.

\bibitem{Rousseau} G.~Rousseau, \emph{Completeness in finite algebras
with a single operation}, Proc.\ Amer.\ Math.\ Soc.\ 18 (1967),
1009--1013.

\bibitem{Sheffer} H.~M.~Sheffer, \emph{A set of five independent
postulates for Boolean algebras, with application to logical
constants}, Trans.\ Amer.\ Math.\ Soc.\ 14 (1913), 481--488.

\bibitem{Sierpinski} W.~Sierpi\'nski, \emph{Sur les fonctions de
plusieurs variables}, Fund.\ Math.\ 33 (1945), 169--173.

\bibitem{Slupecki} J.~S{\l}upecki, \emph{Kryterium pe{\l}no\'sci
wielowarto\'sciowych system\'ow logiki zda\'n}, C.~R.\ S\'eances
Soc.\ Sci.\ Lettres Varsovie Cl.~III 32 (1939), 102--109; English
translation: \emph{A criterion of fullness of many-valued systems of
propositional logic}, Studia Logica 30 (1972), 153--157.

\bibitem{Szendrei} \'A.~Szendrei, \emph{Clones in Universal Algebra},
S\'eminaire de Math\'ematiques Sup\'erieures 99, Presses de
l'Universit\'e de Montr\'eal, 1986.

\bibitem{Webb} D.~L.~Webb, \emph{Generation of any $n$-valued logic
by one binary operation}, Proc.\ Nat.\ Acad.\ Sci.\ USA 21 (1935),
252--254.

\bibitem{Werner} H.~Werner, \emph{Discriminator-Algebras: Algebraic
Representation and Model Theoretic Properties}, Studien zur Algebra
und ihre Anwendungen 6, Akademie-Verlag, Berlin, 1978.

\end{thebibliography}

\appendix

\section{The formal statement layer}\label{app:lean}

For the record, the full definitional trust base of the formal
verification ($R$ an arbitrary commutative ring; \texttt{X 0},
\texttt{X 1} the two variables of a bivariate polynomial ring;
\texttt{bind1 ![a,b] p} denotes the substitution $p(a,b)$; the formal
development uses \texttt{complete} in declaration names for the
generation property):

\begin{verbatim}
inductive Clo (p : MvPolynomial (Fin 2) R) :
    MvPolynomial (Fin 2) R -> Prop
  | atomX : Clo p (X 0)
  | atomY : Clo p (X 1)
  | const (c : R) : Clo p (C c)
  | sub {a b : MvPolynomial (Fin 2) R} :
      Clo p a -> Clo p b -> Clo p (bind1 ![a, b] p)

theorem F2_master (q : MvPolynomial (Fin 2) (ZMod 2)) :
    ~ Clo q (X 0 + X 1) \/ ~ Clo q (X 0 * X 1)

theorem int_master (p : MvPolynomial (Fin 2) Int) :
    ~ Clo p (X 0 + X 1) \/ ~ Clo p (X 0 * X 1)

theorem qOp_complete {R} [CommRing R] [Invertible (2 : R)] :
    forall g, Clo (X 0 ^ 2 - X 1 : MvPolynomial (Fin 2) R) g

theorem complete_iff_two_isUnit (R) [CommRing R] :
    (exists p : MvPolynomial (Fin 2) R, forall g, Clo p g)
      <-> IsUnit (2 : R)
\end{verbatim}

\end{document}
